\chapter{The Ambient Candidate, Boundaries, and the Research Program}
\label{ch12:chapter}

\section{From the finite factor to the ambient package}
\label{ch12:sec:ambient}

The finite-context theory of Chapters~7, 9 and~10 is complete on its own
terms, but it is a theory of the context factor $L$ alone. The Noetic Topos
2.0 candidate is a construction over the \emph{product} of the registered
legacy category with that factor, and the way the candidate is packaged is
itself a mathematical decision with a theorem behind it. This closing
chapter assembles the package, states every boundary that fences it, records
the backward-compatibility conjecture in its exact conjectural form, sets
out the expressivity obligation with its acceptance criteria, situates the
work against the literature, and lists the admission gates that stand
between the present artifact set and any change of status.

\begin{construction}[The successor package]
\label{ch12:con:package}
Let $\cC$ denote the registered legacy category $\mathrm{Noetica}$, assumed
small in the chosen universe, with its proven jointly-surjective topology
$J_{\mathrm{obs}}$ (the legacy topology, written $J_{\mathrm{legacy}}$ in
the baseline chapters). The successor begins with the ambient presheaf topos
\[
\Eamb = \PSh(\cC \times L),
\]
and generates three topology refinements on it:
\begin{enumerate}
\item legacy observational covers in the $\cC$ coordinate;
\item $K_A$ covers in the $L$ coordinate; and
\item $K_R$ covers in the $L$ coordinate.
\end{enumerate}
Writing sheafification-followed-by-inclusion for these refinements as
$M_{\mathrm{obs}}$, $M_A$, and $M_R$, the fully resolved topology and
subtopos are
\[
J_{\mathrm{res}} = J_{\mathrm{obs}} \vee J_A \vee J_R,
\qquad
\Eres = \Sh(\cC \times L,\, J_{\mathrm{res}}).
\]
The underlying topos candidate and the full research object are
\[
\begin{gathered}
\mathrm{NoeTop}_2 := \Eamb,\\
\mathrm{NoeSystem}_2 = \bigl( \mathrm{NoeTop}_2,\; M_{\mathrm{obs}},\;
M_A,\; M_R,\; \Eres,\; \textsc{external mark record} \bigr).
\end{gathered}
\]
$\mathrm{NoeSystem}_2$ is a structured package --- a tuple of a topos, three
modalities, a subtopos, and an external audit layer. It is \emph{not} itself
a topos or a category, and the provenance record is external audit data,
not an internal topos object.
\end{construction}

The definitional choice $\mathrm{NoeTop}_2 := \Eamb$ --- the ambient topos,
not the resolved subtopos --- is forced by the resolved-identity theorem of
Chapter~10: every sheaf for the joined context coverage is separately a
sheaf for each factor topology, so on $\Eres$ both $M_A$ and $M_R$ restrict
to the identity up to canonical isomorphism. Had the successor been defined
as the already-joined sheaf category, the ACBE and RPM reflectors would be
identity on every admitted object and the modal structure would be erased a
second time --- a repetition, one level up, of exactly the collapse the
successor exists to repair. The intended discipline is therefore: raw and
partially resolved evidence lives in the ambient topos, acted on by the
modalities; only jointly resolved objects enter $\Eres$; and no text may
describe $M_A$ or $M_R$ as nontrivial endomodalities of the resolved
subtopos.

\begin{boundary}[The executable schema is narrower than the notation]
\label{ch12:bnd:schema}
Construction~\ref{ch12:con:package} is a research specification, not a
completed implementation. The v0.1 executable schema is deliberately
narrower: it represents ACBE/RPM context-factor families and their joint
refinement only. \emph{Typed observational covers} in the legacy $\cC$
coordinate remain \textsc{open}: no typed construction of observational
cover records, general arrows of $\cC \times L$, full product-site covers,
or their serialization exists in v0.1. Consequently $M_{\mathrm{obs}}$,
$J_{\mathrm{res}}$, and $\Eres$ are, at present, notation with stated
intent, awaiting the typed constructions and proofs listed in the admission
gates below. Relatedly, the once-proposed lift of the distinctness theorem
to $\PSh(\cC \times L)$ by external product of each witness with a
representable $y\cC_0$ is \emph{not} sound as stated --- at the null cover
over $(c, \varnothing)$ the product value need not be a singleton, and
duplicated representable coordinates can obstruct gluing --- so ambient
distinctness of $M_A$ and $M_R$ also remains \textsc{open}; the proved
distinctness result is scoped to $\PSh(L, \Set)$.
\end{boundary}

\begin{remark}[Anatomy of the failed ambient lift, and the candidate
repair]
\label{ch12:rem:ambientlift}
The failed lift is instructive enough to dissect. The proposal was to
carry the witness $F_A$ to the product site as an external product
$y\cC_0 \times F_A$ --- representable in the legacy coordinate, witness
in the context coordinate --- and argue that at $(\cC_0, \Omega)$ the
four-versus-one counting argument survives. Two things go wrong. At a
null cover over a pair $(c, \varnothing)$, the sheaf condition for the
context-coordinate refinement requires a one-element value, but the
product value there has as many elements as $y\cC_0(c)$, which can be
empty or large; the external product does not satisfy the empty-context
discipline that Chapter~10's terminal-empty-value lemma forces on every
sheaf. And along legacy arrows duplicating representable coordinates,
matching families can fail to glue for reasons that have nothing to do
with the context factor, contaminating the witness argument. The
candidate repair replaces $y\cC_0$ by the terminal $\cC$-presheaf,
trivializing the legacy coordinate --- but that repair is only
well-posed after the lifted topology is actually defined, requires
$\cC$ nonempty, and needs naturality checked in both coordinates. None
of that has been done, and the lesson already banked is
methodological: ambient statements do not follow from finite-context
statements by decorating witnesses, and every future lift claim must be
proved against the typed product-site definitions of admission gate 3.
\end{remark}

\begin{remark}[Coordinates, refinements, and the intended evidence flow]
\label{ch12:rem:flow}
Some unpacking of Construction~\ref{ch12:con:package} is useful even
while its product-site layer remains open. The three refinements are
\emph{coordinate-local} by design: an observational cover touches only
the $\cC$ coordinate (it is a legacy jointly-surjective family, paired
with identities on the context coordinate), while an ACBE or RPM cover
touches only the $L$ coordinate (a $K_A$ or $K_R$ family paired with
identities on the legacy coordinate). Generating a topology on the
product from coordinate-local families is standard in outline, but the
candidate treats none of it as automatic: which sieves the joint
generation produces at mixed objects, whether sheafification for a
context-coordinate refinement can be computed objectwise in $\cC$
(pointwise application of Chapter~10's reflector formula), and how the
three refinements interact under joins are precisely the contents of
admission gates 3 and 4 below. Subject to those gates, the intended
evidence flow through the package is a pipeline with typed stages: raw
evidence enters as a presheaf on $\cC \times L$ (an object of $\Eamb$),
carrying arbitrary context-coordinate structure and its marked cover
records; partial resolution applies $M_A$ or $M_R$ --- meaningfully, in
either order, by the composition theorem --- while audit code reads the
marks that the reflectors provably do not see; and joint resolution
lands the object in the modally inert resolved zone, after which only
observational structure remains active. The stage boundaries are
mathematically enforced: Chapter~10's resolved-identity theorem makes
``resolved but still ACBE-active'' a contradiction in terms, which is
exactly the kind of design error the package's typing is meant to make
unwritable.
\end{remark}

\begin{remark}[What the package buys, conditionally]
\label{ch12:rem:buys}
Even in its partly-open state, the package fixes three commitments that
constrain all future work. First, the carrier of modal structure is ambient:
any future theorem about modal dynamics must be stated on $\Eamb$, and any
statement on $\Eres$ must be about resolved, modally inert objects. Second,
provenance is external: the marks travel beside the topology, in
hash-chained records with fail-closed resolution (Chapter~8; executable
checks in Chapter~11), because the
topology provably cannot carry them. Third, the package is honest about its
own type: $\mathrm{NoeSystem}_2$ is a tuple, and calling it a topos would
overclaim structure that only its first component has.
\end{remark}

\section{Gluing boundaries: what Set-valued descent does not license}
\label{ch12:sec:gluing}

The finite theory is a theory of $\Set$-valued descent: gluing means
bijectivity of the canonical map from global sections to matching families
of local sections. Two adjacent temptations --- reading probabilistic or
linear-algebraic consequences out of $\Set$-valued sheaf conditions --- must
be blocked explicitly, because both are natural-sounding and both are
wrong.

\subsection{The probabilistic boundary and the marginal problem}
\label{ch12:sec:probability}

Probability distributions do not inherit $\Set$-valued gluing by analogy.
The obstruction is classical: a family of pairwise-compatible marginal
distributions need not extend to any global joint distribution. This is the
\emph{marginal problem}, and nonextendable compatible families were
exhibited already by Vorob'ev \cite{vorobev}; the sheaf-theoretic
reformulation of such obstructions, in which empirical models are
compatible families over a cover of measurement contexts and contextuality
is precisely the failure of a global section, is developed by Abramsky and
Brandenburger \cite{abramskybrandenburger}. The relevance to this program
is direct and cautionary: the entire motivation for a context site is to
manage heterogeneous evidence, evidence in practice is often probabilistic,
and it is therefore tempting to treat a passing $\Set$-valued sheaf test as
evidence that probabilistic summaries over the same cover cohere. They need
not.

\begin{example}[The obstruction in miniature]
\label{ch12:ex:marginal}
The phenomenon fits in three binary variables. Let $X$, $Y$, $Z$ each
take values in $\{0, 1\}$, and prescribe the three pairwise marginals to
be perfect anticorrelation: each of $(X, Y)$, $(Y, Z)$, $(X, Z)$ is
uniform on $\{(0,1), (1,0)\}$. Every pair of these marginals is
compatible --- each single-variable marginal is uniform on $\{0,1\}$
however computed, so the family agrees on all overlaps. Yet no joint
distribution has all three as marginals: a joint supported on triples
$(x, y, z)$ with $x \neq y$, $y \neq z$, and $x \neq z$ would need three
pairwise-distinct values from a two-element set, so its support is
empty, contradicting total mass one. Pairwise compatibility --- the
exact probabilistic analogue of a matching family --- holds; the global
section does not exist. In gate terms
(Protocol~\ref{ch12:prot:probgates}): the cover
$\{XY, YZ, XZ\}$ is cyclic, so gate 1 does not apply; the gate-2 linear
program for this instance is infeasible, with the counting argument
just given convertible into a Farkas witness; and the family has a
nonvanishing obstruction in the sense of the contextuality literature
\cite{abramskycohomology,abramskyparadox}. Any evidence-aggregation design over the
Noetic context factor that reports probabilistic summaries per cover
piece must therefore treat joint consistency as a per-instance
obligation, not a structural gift of the site.
\end{example}

A probabilistic successor to the present construction must therefore pass
through one of four explicit gates, each of which supplies what the analogy
omits:

\begin{protocol}[The four probabilistic gates]
\label{ch12:prot:probgates}
\begin{enumerate}
\item \textbf{Junction-tree extension.} Restrict to covers whose
intersection structure is acyclic (a junction tree), and supply a proved
extension theorem that compatible marginals over such a cover always admit
a global joint. Acyclicity is exactly the classical regime in which the
marginal problem is solvable, and the proof obligation is the extension
theorem itself, stated for the specific cover class used.
\item \textbf{LP feasibility certificate.} For a fixed finite cover and
finite outcome spaces, existence of a global joint is a finite linear
program: variables the putative joint probabilities, constraints the
marginalization equations and nonnegativity. A probabilistic consistency
claim is then licensed by a feasibility certificate for that LP --- an
explicit feasible point or a solver artifact --- and refuted by an
infeasibility certificate (a Farkas witness). The certificate is
per-instance: it licenses the instance, not the method.
\item \textbf{Obstruction calculation with stated coefficients.} Compute a
cohomological obstruction to extending the compatible family, in the style
of the sheaf-theoretic contextuality literature
\cite{abramskycohomology}, with the coefficient system stated
explicitly. Vanishing of the obstruction in the stated coefficients is the
licensed claim --- no more; such obstructions can be one-directional, and
the direction proved must be the direction cited.
\item \textbf{Quantitative relaxation.} Where exact extension fails or is
not decided, a quantitative claim is permitted only with the relaxation
stated (e.g. nearest extendable family in a named metric), the residual
reported, and the uncertainty of the residual reported. The
sheaf-theoretic literature standardizes one such relaxation: the
contextual fraction of \cite{abramskyfraction}, computed by the same
linear program as gate~2.
\end{enumerate}
No passing $\Set$-valued sheaf test licenses a probabilistic consistency
claim; only a gate artifact of one of these four kinds does.
\end{protocol}

\subsection{The linear and cellular boundary}
\label{ch12:sec:linear}

A second analogy must be blocked in the same motion. The applied-topology
literature attaches spectral invariants --- sheaf Laplacians, harmonic
sections, diffusion dynamics --- to \emph{cellular sheaves} of vector
spaces, and these tools are attractive for evidence aggregation. But a
linear or cellular sheaf is additional structure, not a reinterpretation:
it requires finite-dimensional stalks (or controlled Hilbert-space stalks),
chosen inner products, and weighted coboundary or restriction data before a
Laplacian is even defined \cite{hansenghrist}. Nothing in the present
construction creates that structure. The context lattice $L$ carries
$\Set$-valued presheaves; no stalk is a vector space, no inner product has
been chosen, and no coboundary operator exists. Any future spectral or
diffusion-style claim over the Noetic context factor requires a separate
linear representation with its own registered definitions and proofs, and
bounded Hilbert-valued cellular sheaves need not be finite-dimensional, so
even the passage from finite to infinite dimensions requires care rather
than analogy.

To make the missing structure concrete: a cellular sheaf in the sense of
\cite{hansenghrist} assigns a vector space to each cell of a complex and
a linear map to each incidence, and its Laplacian is built from a chosen
coboundary operator and inner products on stalks --- so that
``harmonic'' sections, spectral gaps, and consistency energies are all
statements about \emph{that chosen linear data}, and change when the
inner products change. The Noetic context factor supplies only the
incidence skeleton (the inclusion order on $L$) and $\Set$-valued
stalks. Between here and any Laplacian-style consistency score lie, at
minimum: a linearization of section sets (itself a modeling decision
with alternatives), a choice of inner products with a stated rationale,
a weighting of the coboundary reflecting cover semantics, and proofs
that the resulting invariants respect the reflectors of Chapter~10 or
an argument for why they need not. Each is a registered-definition-sized
gap, none is closed by this monograph, and any program text that
computes a ``coherence energy'' over the context factor without them is
operating outside the licensed mathematics.

\section{Backward compatibility as a conjectural comparison}
\label{ch12:sec:compat}

The successor would be of limited use if it stranded the legacy theory. The
intended relationship is a comparison, and its status is exactly
\textsc{conjecture} with \textsc{proof required} --- the only record in the
finite theorem set that is not backed by a proof, and the monograph states
it in conjectural form for that reason.

\begin{conjecture}[Product-site round trip; ledger
\textsc{thm-nt2-product-site-backward-compatibility}, \textsc{proof
required}]
\label{ch12:conj:roundtrip}
Let $\Eleg = \Sh(\cC, J_{\mathrm{legacy}})$ be the legacy sheaf topos,
preserved verbatim. After atomic resolution of the context factor there is
an equivalence
\[
\Eres \;\simeq\; \Eleg^{4},
\]
under which, for a registry-fixed neutral atom $\omega_{00}$, the embedding
\[
\iota : X \longmapsto (X, 1, 1, 1)
\]
--- the legacy object in the neutral channel, terminal objects in the
remaining three channels --- is fully faithful, and evaluation at the
neutral channel satisfies a specified natural isomorphism
\[
\mathrm{ev}_{00} \circ \iota \;\cong\; \mathrm{Id}_{\Eleg}.
\]
\end{conjecture}

\begin{remark}[Why the shape is plausible and why that is not a proof]
\label{ch12:rem:whyplausible}
The finite-context theory makes the \emph{shape} of the conjecture
plausible: the joined context coverage resolves the four singleton atoms
(Chapter~10), and sheaves on a discrete four-point context factor should
decompose into four independent components. But the conjecture is a
statement about the \emph{product} site $\cC \times L$ with the
observational topology in play, and none of the required ingredients ---
the typed product-site covers (Boundary~\ref{ch12:bnd:schema}), the
comparison functors, the full-faithfulness argument, the coherence of the
equivalence with $J_{\mathrm{obs}}$ --- currently exists. A proof would
need, at minimum: a typed construction of $J_{\mathrm{res}}$ on
$\cC \times L$; a computation showing sheafification for the joined context
coordinate is objectwise in $\cC$ (the product-site lift of the composition
theorem, itself open); and the equivalence with naturality in all data.
\end{remark}

\begin{boundary}[Two round trips, and neither is complete]
\label{ch12:bnd:tworoundtrips}
The compatibility contract requires \emph{two distinct} round trips, and
they must not be conflated:
\begin{enumerate}
\item \textbf{The categorical round trip.} The equivalence and natural
isomorphism of Conjecture~\ref{ch12:conj:roundtrip}. This is mathematics,
and its status is \textsc{proof required}.
\item \textbf{The serialization round trip.} The implementation requirement
that registered identifiers and canonical payload hashes survive migration
into and out of the new context factor \emph{by strict equality}. This is
engineering with a cryptographic acceptance test, and it is likewise
incomplete: the v0.1 context-factor checker does not establish it.
\end{enumerate}
A categorical natural isomorphism is not a bit-for-bit serialization
receipt, and a passing serialization test would not prove the equivalence.
Both are required; possession of one licenses nothing about the other.
Additionally, because subcanonicity of the legacy site is not derivable
from the registered axioms (the baseline countermodel), bare legacy site
objects are represented in the sheaf topos by their sheafifications
$a_{\mathrm{legacy}}(yX)$ before any marking or migration --- never by
silently treating $yX$ as a sheaf \cite{stacks00WO}. This monograph claims
no backward compatibility; no legacy theorem is promoted merely because it
has been serialized into the new context factor.
\end{boundary}

\begin{remark}[What full faithfulness would and would not buy]
\label{ch12:rem:fullfaithful}
We spell out the operational content the conjecture would have if
proved, because that content is the reason the conjecture is
registered at all. Full faithfulness of $\iota$ would say that legacy
objects and \emph{all their morphisms} survive the move into the
neutral channel without collapse or spurious identification: maps
between embedded legacy objects in the resolved successor would
correspond exactly to legacy maps. The natural isomorphism
$\mathrm{ev}_{00} \circ \iota \cong \mathrm{Id}$ would say the move is
reversible at the object level by evaluation. Together they would
license the migration discipline the program wants: analyze in the
successor, project back, and lose nothing that the legacy semantics
could see. What even a full proof would \emph{not} buy: preservation of
anything the legacy semantics could not see (marks are external on both
sides and migrate by their own serialization contract); any statement
about the three non-neutral channels, which the embedding deliberately
fills with terminal objects; and any transfer of legacy theorem
\emph{statuses} --- a legacy theorem cited in successor context still
carries its own ledger record, including the legacy subcanonicity gap.
The neutral channel being registry-fixed matters for the same
discipline: if different artifacts chose different neutral atoms, their
embeddings would disagree by a nontrivial symmetry of the context
factor, and cross-artifact comparisons would silently break; fixing
$\omega_{00}$ in the registry makes the choice auditable.
\end{remark}

\section{The expressivity boundary and the cross-context obligation}
\label{ch12:sec:expressivity}

Suppose Conjecture~\ref{ch12:conj:roundtrip} were proved tomorrow. A
further boundary would immediately bind, and it is in some ways the deepest
one in the monograph: $\Eleg^4$ is four \emph{independent} legacy channels.
The discrete resolved context factor contains no remaining overlap, no
transition law, and no cross-atom coupling. The composition theorem lands
every composite in the atomic subcategory, where the modalities freeze
(Chapter~10, identity-on-atomic corollary); after resolution there is
nothing left for the context factor to say about how the four channels
relate. The sixteen-context construction is therefore a finite
\emph{separation and provenance reference model} --- it proves that
distinct, nontrivial, commuting context modalities are possible where the
legacy clauses provably could not provide them --- and not a substantively
rich final semantics for contextual cognition.

The registered proof-obligation ledger carries this as a first-class
record, and its content is reproduced here in full because it, rather than
any theorem, determines the shape of the next stage of the program. The
\emph{ledger} is not, and the distinction matters to a reader auditing
completeness: what follows is one obligation record, not the register. How many records the ledger holds in total is not among the
figures registered in the source artifacts of Chapter~1,
Section~1.\ref{ch01:sec:sources}, so under Rule~\ref{ch01:rule:claims}(iv)
that count is declared unspecified here rather than guessed; the register
itself lives in the proof-obligation ledger v0.1, cited there with its
digest. What this chapter does reproduce whole is the ten admission gates
of Section~\ref{ch12:sec:gates}.

\begin{invariant}[Obligation \textsc{obl-nt2-cross-context-semantics},
v0.1]
\label{ch12:inv:obligation}
\emph{Kind:} design-and-proof obligation. \emph{Statement:} choose and
formalize one conservative successor contract: either (a) the structured
ambient modal package is the principal Noetic Topos 2.0 semantics, or (b)
the resolved semantics gains typed overlap or transition data with proved
coherence. The finite discrete product $\Eleg^4$ alone supplies separation
and provenance channels but no cross-atom coupling. \emph{Assumptions:} the
joined context coverage resolves the four singleton atoms; the proposed
product-site comparison is $\Eres \simeq \Eleg^4$. \emph{Dependencies:} the
join theorem (ledger \textsc{thm-nt2-join-discrete-context}) and the
backward-compatibility conjecture (ledger
\textsc{thm-nt2-product-site-backward-compatibility}). \emph{Review state:}
\textsc{unassessed}; \emph{resolution status:} \textsc{open}; \emph{ring:}
4.

\emph{Acceptance criteria} (all four required):
\begin{enumerate}
\item a single conservative successor contract is selected without
asserting unproved equivalences;
\item the selected contract has typed definitions and explicit coherence
obligations;
\item any claimed coherence theorem has a schema-valid theorem-ledger
record and a reviewable proof artifact; and
\item the four-independent-channel limitation remains explicit unless
cross-context coupling is proved.
\end{enumerate}
\emph{Claim boundary of the record itself:} the current finite model is a
reference construction; it is not licensed as the final substantive Noetic
Topos 2.0 semantics.
\end{invariant}

\begin{remark}[The two branches, concretely]
\label{ch12:rem:branches}
Branch (a) names the ambient package --- topos plus modalities plus
resolved subtopos plus marks --- as the semantics, accepting that the
interesting structure lives in the \emph{unresolved} region and that
resolution is an endpoint, not a home. Branch (b) enriches the resolved
side: typed overlap objects, transition data between atoms, or a
coarser-than-discrete resolved context structure, each with a coherence
theorem. The branches are not exclusive forever, but the obligation demands
that \emph{one} be selected and formalized first, precisely to prevent the
program from oscillating between an unformalized version of each. v0.1
licenses neither branch, and this monograph selects neither; selection is a
gated decision with its own review.
\end{remark}

\section{Related work and novelty status}
\label{ch12:sec:related}

The mathematical background this monograph stands on is entirely standard,
and no novelty is claimed for any of it: sites, sieves, and the passage
from a coverage to a Grothendieck topology; sheafification as a left-exact
left adjoint and its universal property; representable sheaves and
subcanonicity; and the Grothendieck-topos framework generally
\cite{maclane, stacks00ZC, stacks00W1, stacks00WO}; with the local-operator
perspective on modalities available in \cite{johnstone}. Two applied
literatures motivate the explicit boundaries of
Section~\ref{ch12:sec:gluing} without being incorporated into the
construction: the sheaf-theoretic treatment of nonlocality and
contextuality, where global-section failure over a cover of measurement
contexts is the formal content of contextuality
\cite{abramskybrandenburger}, with the underlying marginal-problem
obstructions classical \cite{vorobev}; and the spectral theory of cellular
sheaves \cite{hansenghrist}, which shows exactly how much added linear
structure a Laplacian requires and therefore exactly what the present
$\Set$-valued theory lacks.

Against that background, the candidate's specific contributions are three:
the collapse diagnosis as applied to the registered legacy ACBE/RPM
clauses; the transverse two-partition finite context factor with its exact
reflector calculus (Chapter~10); and the parallel hash-chained provenance
layer with fail-closed resolution (Chapter~8, corroborated in
Chapter~11). Their status must be
stated exactly: all theorem-ledger records carry
\textsc{novelty unassessed}. No dedicated literature search has been
performed for the relationship of the marked-cover construction to colored
sites, indexed categories, or stacks; nor for whether the explicit
product-over-blocks reflector for partition coverages on a powerset
lattice appears elsewhere (it is elementary enough that it plausibly
does); nor for prior uses of transverse partition pairs to manufacture
commuting distinct modalities. This monograph accordingly does not call
any of its results new theorems in the scholarly sense. The claim made is
weaker and auditable: the results are \emph{internally novel} to the
program, correctly scoped, and fully proved or fully flagged.

\section{The admission gates}
\label{ch12:sec:gates}

Noetic Topos 2.0 may leave Ring~4 only after all ten of the following
close. The list is registered; this section adds, for each gate, what
closing it concretely requires.

\begin{protocol}[The ten admission gates]
\label{ch12:prot:gates}
\begin{enumerate}
\item \textbf{Independent category-theory review} of the collapse and
countermodel propositions and of the finite-context theorems through the
composition theorem. Closing requires a named, independent reviewer;
internal re-derivation, however careful, does not qualify.
\item \textbf{Proof-assistant formalization} of the finite site, the
generated topologies, the witness presheaves, the reflector universal
properties, the join, and the natural isomorphisms --- the full
sheaf-theoretic layer excluded from the current Lean scope (Chapter~11).
\item \textbf{Typed observational and product-site covers:} a typed
construction of observational covers and general arrows in
$\cC \times L$, discharging Boundary~\ref{ch12:bnd:schema}.
\item \textbf{The product-site lift} of the composition theorem: its
coherence in the $\cC$ coordinate and its interaction with
$J_{\mathrm{obs}}$. Until this closes, commutation is a finite-context
fact only.
\item \textbf{The backward-compatibility conjecture}
(Conjecture~\ref{ch12:conj:roundtrip}) with \emph{both} round trips of
Boundary~\ref{ch12:bnd:tworoundtrips}: the categorical equivalence with
proofs, and the strict-equality serialization round trip with tests.
\item \textbf{A canonical cross-context decision} under obligation
\textsc{obl-nt2-cross-context-semantics}
(Invariant~\ref{ch12:inv:obligation}), meeting all four acceptance
criteria.
\item \textbf{Probabilistic descent}: separate probability-gluing theorems
or feasibility certificates through the four gates of
Protocol~\ref{ch12:prot:probgates}, for any probabilistic use of the site.
\item \textbf{A dedicated novelty search} for the marked/provenance
construction and the reflector calculus, converting
\textsc{novelty unassessed} into an assessed status in either direction.
\item \textbf{Registry reconciliation and IP review}: reconciliation with
the shared symbol and theorem registries, and authorship, disclosure,
citation, and IP review of the whole artifact set.
\item \textbf{External referee review}, with explicit confirmation that no
empirical or product wording was licensed anywhere in the artifact set.
\end{enumerate}
\end{protocol}

Gates 1, 2, and 10 concern trust in the mathematics; gates 3, 4, and 5
concern the ambient and legacy extensions; gates 6 and 7 concern semantics
beyond the finite reference model; gates 8 and 9 concern scholarly and
institutional hygiene. They are conjunctive: no subset suffices, and no
gate is closed by the drafting of this monograph.

The gates also carry an implicit dependency order that any realistic
program plan must respect. Gate 3 (typed product-site covers) is
upstream of gates 4 and 5: neither the product-site lift nor the
backward-compatibility equivalence can even be \emph{stated} precisely
until the typed constructions exist, which is why both are conjectural
or open rather than merely unproved. Gate 5 is upstream of gate 6: the
cross-context obligation's assumptions include the proposed comparison
$\Eres \simeq \Eleg^4$, so a canonical semantics decision made before
the comparison is settled would rest on an unproved premise ---
permitted only under branch (a), which does not need the comparison, and
then only with the four-channel limitation stated as conditional. Gates
1 and 2 can proceed in parallel with everything else and de-risk it:
independent review or formalization that found an error in the
finite-context theory would invalidate downstream work cheaply rather
than expensively. Gates 8, 9, and 10 are terminal: novelty search,
registry and IP review, and external refereeing evaluate the finished
artifact set, and doing them against a moving target would waste their
one-shot evidential value. None of this ordering relaxes conjunctivity;
it only says in which order the program can stop being wrong soonest.

\section{Conclusion}
\label{ch12:sec:conclusion}

The monograph's argument closes where it began, with the two collapse
mechanisms. The legacy ACBE and RPM cover clauses could not carry modal
structure because identities entered their sieves: the reflexive descent
map directly, the RPM counit through its section by sieve closure. That is
not a defect of topos machinery but a theorem about what a Grothendieck
topology can remember, and it dictated both halves of the repair. The
topological half changes the generators: a finite context lattice with two
transverse partition coverages whose global pieces are proper, so that the
generated sieves are nonmaximal and the modalities live. On that factor
the theory is complete and exact --- genuine subcanonical topologies;
distinct, non-identity modalities separated by explicit witnesses; a
computed join; explicit reflectors verified directly against the universal
property; and the composition theorem, which is corrective as well as
constructive: the modalities are distinct but \emph{commuting}, either
composite is sheafification for the join, and no artifact licenses a
noncommutation claim. The provenance half moves what the topology forgets
into hash-chained marked records that fail closed, checked at the
serialization level by the same certificate that checks the theorems.

Around that completed core, every remaining surface is fenced. The ambient
package is named honestly as a structured tuple whose executable schema is
narrower than its notation; the probabilistic and linear analogies are
blocked by explicit gates; backward compatibility is a conjecture with two
incomplete round trips; the resolved semantics is four independent
channels until coupling is proved, under a registered obligation with
acceptance criteria; novelty is unassessed; and ten conjunctive gates
stand between the artifact set and any exit from Ring~4. The monograph's
contribution is the pairing: a small theory proved to the bottom, and a
boundary discipline maintained with the same precision as the proofs.
Preserving that pairing --- promoting nothing that the ledger does not
promote, while defending everything it does --- is the standard against
which any successor to this successor should be judged.
